N≫1 个彼此无相互作用、可区分的二能级单位，单粒子能级分别为 0 与 ε>0。固定激发粒子数 n、总内能 U=nε。考虑宏观热力学极限并用 Boltzmann 组合熵 S=kB ln[binom(N,n)]。
(1) 用 Stirling 近似求激发比例 x=n/N 下的熵；(2) 由 1/T=(∂S/∂U)N 得到温度与 x 的关系，识别负温度区域；(3) 求正则体系的激发概率和热容，验证正、负温度延拓的一致性；(4) 说明 x→1/2、x→1⁻ 极限及为什么这种负温度不是“比零更冷”。
